\documentclass{article}
\usepackage[margin=1in]{geometry}
\usepackage{tikz}
\usepackage{amsmath}
\usepackage{amssymb}
\usepackage{multicol}
\usepackage{import}

\import{../}{labflowchart.sty}

\begin{document}
\title{SCH4U Lab: pH of Salts and Ka}
\author{Aaron Huang \\
Lab Partners: Maya Nasr, Arshan Shahkarami, Ryan Huang}
\date{April 28, 2026}
\maketitle

% \section*{Purpose}
% \begin{enumerate}
%     \item To measure the pH of salt solutions
%     \item To determine the Ka of weak acids using pH measurements
%     \item To prepare and test a buffer solution
% \end{enumerate}

\section*{Introduction}

This lab explores the pH of salt solutions, the acid dissociation of weak acids, and the properties of buffer solutions. When a neutralization occurs, a salt is formed. Neutral salts form when a strong acid and a strong base react. A non-neutral salt forms when a weak acid or a weak base is involved, it exhibits either acidic or basic properties due to hydrolysis reactions in water. 

Weak acids/bases only partially dissociate in aqueous solution. The strength of a weak acid or bases is quantified by its acid dissociation constant, $K_a$ or $K_b$ respectively. By measuring the pH of buffer solutions containing a weak acid/weak base and its conjugate base/acid, we can determine the $K_a$/$K_b$ value. A buffer solution resists changes in pH when small amounts of acid or base are added, making it useful in many chemical and biological applications.

This lab consists of three parts. First, the pH of different salt solutions will be measured using a pH probe. Second, buffer solutions will be prepared by combining a weak acid with its conjugate base, and their pH will be recorded to calculate Ka values. Finally, the buffer's ability to resist pH changes will be tested by adding a strong base and comparing the pH change to that of pure water treated the same way.

If solutions contact skin, rinse immediately as they may cause skin irritation and handle glassware with care.

\section*{Procedure}

\subsection*{Part 1: pH of Salt Solutions}

\begin{procedureflowchart}
    \lstep{Beaker (100 mL)}
    \rstep{Prepare clean beaker}
    
    \lstep{Graduated Cylinder}
    \rstep{Measure and add 20.0 mL of 0.100 M NaCl}
    
    \lstep{pH Probe}
    \rstep{Calibrate probe and measure pH}
    \rstep{Record pH}
    
    \lstep{Deionized Water}
    
    \rstep{Repeat for Na$_2$CO$_3$ and NaHSO$_3$}
\end{procedureflowchart}

\subsection*{Part 2: Ka of a Weak Acid}

\begin{procedureflowchart}
    \lstep{Beaker (250 mL)}
    \rstep{Prepare clean beaker}
    
    \lstep{Graduated Cylinder}
    \rstep{Measure and add 20.0 mL of 0.1 M NaHSO$_3$}
    \rstep{Measure and add 20.0 mL of 0.100 M Na$_2$SO$_3$}
    
    \lstep{Glass Rod}
    \rstep{Stir to mix}
    
    \lstep{pH Probe}
    \rstep{Measure and record pH}
    \rstep{Rinse probe}
    
    \rstep{Repeat with CH$_3$COOH and NaCH$_3$COO solutions}
\end{procedureflowchart}

\subsection*{Part 3: Buffer Solution Testing}

\begin{procedureflowchart}
    \lstep{Buffer Solution (from Part 2)}
    \rstep{Prepare 40.0 mL of NaHSO$_3$/Na$_2$SO$_3$ buffer}
    
    \lstep{Graduated Cylinder}
    \rstep{Add 10.0 mL of NaOH to buffer}
    
    \lstep{pH Probe}
    \rstep{Measure and record pH}
    
    \lstep{Beaker}
    \rstep{Prepare 40 mL of deionized water}
    
    \lstep{Graduated Cylinder}
    \rstep{Add 10.0 mL of NaOH to water}
    
    \lstep{pH Probe}
    \rstep{Measure and record pH}
\end{procedureflowchart}

\section*{Observations}


\begin{table}[ht!]
    \centering
    \caption{pH Measurements of Salt Solutions}
    \label{tab:part1-observations}
    \vspace{0.125 in}
    \renewcommand{\arraystretch}{2.5}
    \begin{tabular}{|c|c|}
        \hline
        \textbf{Solution} & \textbf{pH} \\
        \hline
        0.100 M NaCl & 5.6 \\
        \hline
        0.100 M Na$_2$CO$_3$ & 9.9 \\
        \hline
        0.100 M NaHSO$_3$ & 4.3 \\
        \hline
    \end{tabular}
\end{table}


\begin{table}[ht!]
    \centering
    \caption{pH of Buffer Solutions}
    \label{tab:part2-observations}
    \vspace{0.125 in}
    \renewcommand{\arraystretch}{2.5}
    \begin{tabular}{|c|c|c|}
        \hline
        \textbf{Buffer} & \textbf{Acid} & \textbf{pH} \\
        \hline
        NaHSO$_3$/Na$_2$SO$_3$ & HSO$_3^-$ & 6.5 \\
        \hline
        NaCH$_3$COO/CH$_3$COOH & CH$_3$COOH & 6.6 \\
        \hline
    \end{tabular}
\end{table}


\begin{table}[ht!]
    \centering
    \caption{pH Changes Upon Addition of NaOH}
    \label{tab:part3-observations}
    \vspace{0.125 in}
    \renewcommand{\arraystretch}{2.5}
    \begin{tabular}{|c|c|c|c|}
        \hline
        \textbf{Solution} & \textbf{Initial pH} & \textbf{After 10 mL NaOH} & \textbf{pH Change} \\
        \hline
        Buffer (HSO$_3^-$/SO$_3^{2-}$) & 6.5 & 6.9 & 0.4 \\
        \hline
        Deionized Water & 7.0 & 11.0 & 4.0 \\
        \hline
    \end{tabular}
\end{table}
\newpage
\section*{Analysis}

\subsection*{Part 1: Predicted pH of Salt Solutions}

\subsubsection*{Na$_2$CO$_3$ pH Calculation}

\begin{equation*}
\text{CO}_3^{2-} + \text{H}_2\text{O} \rightleftharpoons \text{HCO}_3^- + \text{OH}^-
\end{equation*}

Calculate the $K_b$ value of CO$_3^{2-}$ ($K_{a2} = 4.7 \times 10^{-11}$):

\noindent
\begin{minipage}[t]{0.5\textwidth}
    \begin{align*}
        K_b &= \frac{K_w}{K_{a2}} \\
        &= \frac{1.0 \times 10^{-14}}{4.7 \times 10^{-11}} \\
        &= 2.127659 \times 10^{-4} \\
        &= 2.1 \times 10^{-4}
    \end{align*}
\end{minipage}

\vspace{0.1in}
\begin{table}[h!]
    \centering
    \renewcommand{\arraystretch}{2.0}
    \begin{tabular}{|c|c|c|c|}
        \hline
        & [CO$_3^{2-}$] & [HCO$_3^-$] & [OH$^-$] \\
        \hline
        \textbf{I} (mol/L) & 0.100 & 0 & 0 \\
        \hline
        \textbf{C} (mol/L) & $-x$ & $+x$ & $+x$ \\
        \hline
        \textbf{E} (mol/L) & $0.100-x$ & $x$ & $x$ \\
        \hline
    \end{tabular}
\end{table}

\begin{equation*}
K_b = \frac{[\text{HCO}_3^-][\text{OH}^-]}{[\text{CO}_3^{2-}]} = \frac{x^2}{0.100-x} = 2.127659 \times 10^{-4}
\end{equation*}

\begin{equation*}
x^2 + 2.127659 \times 10^{-4}x - 2.127659 \times 10^{-5} = 0
\end{equation*}

$x = [\text{OH}^-] = 4.50749 \times 10^{-3}$ mol/L $\rightarrow 4.5 \times 10^{-3}$ mol/L

\vspace{0.1in}
\noindent
\begin{minipage}[t]{0.45\textwidth}
    \begin{align*}
        \text{pOH} &= -\log(4.50749 \times 10^{-3}) \\
        &= 2.34606 \\
        &= 2.35
    \end{align*}
\end{minipage}
\hfill
\begin{minipage}[t]{0.45\textwidth}
    \begin{align*}
        \text{pH} &= 14.00 - 2.34606 \\
        &= 11.65394 \\
        &= 11.65
    \end{align*}
\end{minipage}

\vspace{0.1in}
$\therefore$ The predicted pH of Na$_2$CO$_3$ is 11.65 (experimental: 9.9)

\subsubsection*{NaHSO$_3$ pH Calculation}


\noindent
\begin{minipage}[t]{0.45\textwidth}
    \textbf{As an acid:}
    \begin{align*}
        K_{a2} &= 6.2 \times 10^{-8}
    \end{align*}
\end{minipage}
\hfill
\begin{minipage}[t]{0.45\textwidth}
    \textbf{As a base:}
    \begin{align*}
        K_b &= \frac{K_w}{K_{a1}} = \frac{1.0 \times 10^{-14}}{1.3 \times 10^{-2}} \\
        &= 7.692307 \times 10^{-13} \\
        &= 7.7 \times 10^{-13}
    \end{align*}
\end{minipage}

\vspace{0.1in}
Since $K_{a2} > K_b$, HSO$_3^-$ acts as an acid:
\begin{equation*}
\text{HSO}_3^- + \text{H}_2\text{O} \rightleftharpoons \text{SO}_3^{2-} + \text{H}_3\text{O}^+
\end{equation*}

\begin{table}[h!]
    \centering
    \renewcommand{\arraystretch}{2.0}
    \begin{tabular}{|c|c|c|c|}
        \hline
        & [HSO$_3^-$] & [SO$_3^{2-}$] & [H$_3$O$^+$] \\
        \hline
        \textbf{I} (mol/L) & 0.100 & 0 & 0 \\
        \hline
        \textbf{C} (mol/L) & $-x$ & $+x$ & $+x$ \\
        \hline
        \textbf{E} (mol/L) & $0.100-x$ & $x$ & $x$ \\
        \hline
    \end{tabular}
\end{table}

\noindent
\begin{minipage}[t]{0.45\textwidth}
    \textbf{Assume:} \\
    $\frac{0.100}{6.2 \times 10^{-8}} = 1.6129 \times 10^{6} > 1000$ \\
    Assumption is valid.
\end{minipage}
\hfill
\begin{minipage}[t]{0.45\textwidth}
    \textbf{Verify:} \\
    $\frac{0.100 - x}{0.100} \times 100\% = 99.9\% > 95\%$ \\
    Verification successful.
\end{minipage}

\begin{equation*}
K_{a2} = \frac{[\text{SO}_3^{2-}][\text{H}_3\text{O}^+]}{[\text{HSO}_3^-]} \approx \frac{x^2}{0.100} = 6.2 \times 10^{-8}
\end{equation*}

\noindent
\begin{minipage}[t]{0.45\textwidth}
    \begin{align*}
        x^2 &= 6.2 \times 10^{-9} \\
        x &= 7.8740 \times 10^{-5} \\
        &= 7.9 \times 10^{-5}
    \end{align*}
\end{minipage}
\hfill
\begin{minipage}[t]{0.45\textwidth}
    \begin{align*}
        \text{pH} &= -\log(7.8740 \times 10^{-5}) \\
        &= 4.1038 \\
        &= 4.10
    \end{align*}
\end{minipage}

\vspace{0.1in}
$\therefore$ The predicted pH of NaHSO$_3$ is 4.10 (experimental: 4.3)

\subsection*{Part 2: $K_a$ Determinations and Percent Error}

\subsubsection*{$K_a$ of HSO$_3^-$}

\begin{equation*}
\text{pH} = \text{p}K_a + \log\left(\frac{[\text{A}^-]}{[\text{HA}]}\right) = \text{p}K_a + \log(1) = \text{p}K_a
\end{equation*}

\noindent
\begin{minipage}[t]{0.45\textwidth}
    \begin{align*}
        \text{p}K_a &= 6.5 \\
        K_a &= 10^{-6.5} \\
        &= 3.16227 \times 10^{-7} \\
        &= 3.2 \times 10^{-7}
    \end{align*}
\end{minipage}
\hfill
\begin{minipage}[t]{0.45\textwidth}
    \begin{align*}
        \% \text{ error} &= \left| \frac{3.16227 \times 10^{-7} - 6.2 \times 10^{-8}}{6.2 \times 10^{-8}} \right| \times 100\% \\
        &= 409.72\% \\
        &= 4.1 \times 10^{2}\%
    \end{align*}
\end{minipage}

\vspace{0.1in}
$\therefore$ The experimental $K_a$ of HSO$_3^-$ was determined to be $3.2 \times 10^{-7}$ with a $4.1 \times 10^{2}\%$ percent error.

\subsubsection*{$K_a$ of CH$_3$COOH}

\noindent
\begin{minipage}[t]{0.45\textwidth}
    \begin{align*}
        \text{p}K_a &= 6.6 \\
        K_a &= 10^{-6.6} \\
        &= 2.51188 \times 10^{-7} \\
        &= 2.5 \times 10^{-7}
    \end{align*}
\end{minipage}
\hfill
\begin{minipage}[t]{0.45\textwidth}
    \begin{align*}
        \% \text{ error} &= \left| \frac{2.51188 \times 10^{-7} - 1.8 \times 10^{-5}}{1.8 \times 10^{-5}} \right| \times 100\% \\
        &= 98.604\% \\
        &= 98.6\%
    \end{align*}
\end{minipage}

\vspace{0.1in}
$\therefore$ The experimental $K_a$ of CH$_3$COOH was determined to be $2.5 \times 10^{-7}$ with a $98.6\%$ percent error.

% \subsection*{Part 3: Buffer Solution Testing}

% \subsubsection*{Buffer Capacity Analysis}

% Comparison of pH stability between the HSO$_3^-$/SO$_3^{2-}$ buffer and deionized water upon addition of $10.0$ mL of NaOH:

% \noindent
% \begin{minipage}[t]{0.45\textwidth}
%     \textbf{Buffer Solution:}
%     \begin{align*}
%         \Delta\text{pH} &= |6.9 - 6.5| \\
%         &= 0.4
%     \end{align*}
% \end{minipage}
% \hfill
% \begin{minipage}[t]{0.45\textwidth}
%     \textbf{Deionized Water:}
%     \begin{align*}
%         \Delta\text{pH} &= |11.0 - 7.0| \\
%         &= 4.0
%     \end{align*}
% \end{minipage}

% \vspace{0.1in}
% The buffer solution effectively resisted the pH change, demonstrating the capacity of the conjugate acid-base pair to neutralize added ions.
\newpage
\section*{Discussion}
% Discussion Questions:
% 1. Compare the predicted pH of the NaHSO3 solution with the actual experimental result (part 1). Explain any sources of error.
% 2. Provide two explanations for the percent error in the Ka results for the CH3COOH/NaCH3COO equilibrium. (part 2)
% 3. Explain, using Le Chatelier's principle, why the change in pH of the buffer solution was less that the water when adding the NaOH strong base. Be specific. (part 3)

1. The predicted pH of the NaHSO$_3$ solution was 4.104, while the actual experimental result was 4.3. The difference in pH can be attributed to the conditions during the experiment not being SATP, while the calculations were based on SATP conditions. At lower temperatures, the equilibrium of the reaction would shift to the left, which would lead to less dissociation of HSO$_3^-$ and a higher pH. 

2. One possible explanation for the percent error in the Ka of CH$_3$COOH and NaCH$_3$COO is the accuracy of the pH meter. Since the pH probe could only measure pH to one significant digit, the logarithmic relationship between pH and $[H^+]$ would magnify any errors in the pH measurement, leading to a large percent error in the Ka value. 

Another possibility is that the concentration of the CH$_3$COOH and NaCH$_3$COO was not accurately known. Due to the age of the solutions, significant evaporation may have occurred, increasing the concentration of the solutions and thus skewing the pH values. The discrepancy between experimental and theoretical values is likely due to a combination of these factors. 

3. When a strong base, such as NaOH, was added to the HSO$_3^-$/SO$_3^{2-}$ buffer solution listed in Table 3, the $OH^-$ ions reacted with the $H_3O^+$ ions to form water ($OH^- + H_3O^+ \rightarrow 2H_2O$). This significantly reduced the concentration of $H_3O^+$ in the solution. According to Le Chatelier's principle, the system responded to this stress by shifting the equilibrium to the right to produce more $H_3O^+$ ions and restore equilibrium. This replaced most of the $H^+$ ions that were consumed, thus stabilizing the pH. In contrast, when NaOH was added to deionized water, there was no system in place to restore the $H_3O^+$ concentration, causing the pH to increase significantly as shown in Table \ref{tab:part3-observations}.

\section*{Conclusion}
In conclusion, the pH of various salt solutions was measured and predicted, and the experimental $K_a$ values for HSO$_3^-$ and CH$_3$COOH were successfully determined. These findings addressed the initial purpose of evaluating hydrolysis reactions and the buffering capacity of weak acid systems.

\end{document}

